\documentclass[10pt,a4paper]{amsart}
\usepackage[margin=19mm]{geometry}
\usepackage{amsmath,amssymb,mathtools}
\usepackage[scr=rsfs,cal=euler]{mathalfa}
\usepackage{xfrac}
\usepackage[unicode,hidelinks,bookmarks=false]{hyperref}
\newcommand{\ie}{i.e.\ }
\newcommand{\cf}{cf.\ }
\newcommand{\resp}{respectively\ }
\newcommand{\Subsection}{\subsection}
\providecommand{\nobreakdash}{\nobreak\hbox{-}}
\setlength{\emergencystretch}{2em}
\multlinegap0pt
\usepackage{bm}
\usepackage{bbm}
\usepackage{dsfont}
\usepackage{xspace}
\usepackage{cancel}
\usepackage{mathtools}
\usepackage{amsthm}
\makeatletter
\@ifundefined{satz}{\newtheorem{satz}{Satz}}{}
\@ifundefined{proposition}{\newtheorem{proposition}[satz]{Proposition}}{}
\makeatother
\renewcommand{\geq}{\geqslant}
\renewcommand{\leq}{\leqslant}
\title[Optimal Parametrizations and Valuations]{Optimal Parametrizations and Valuations}
\author{Reinhold H\"ubl, Craig Huneke, Sarasij Maitra, Vivek Mukundan}
\date{Source: arXiv:2509.17139v1 (CC BY 4.0)}
\begin{document}
\maketitle

\noindent\emph{Source and licence.} Optimal Parametrizations and Valuations. Version arXiv:2509.17139v1 (CC BY 4.0), distributed under the Creative Commons Attribution 4.0 International licence, as recorded in the arXiv licence metadata for this exact version and shown on the abstract page https://arxiv.org/abs/2509.17139v1. Full rights notice: licenses/arxiv-2509.17139-CC-BY-4.0.txt.\par\smallskip

\noindent\textit{Editorial scope.} Counted unit: the Proposition whose source label is \texttt{prop:MM\_4\_4\_1} in the article (source0.tex lines 783--798, source label \texttt{prop:MM\_4\_4\_1}), delivered together with the article's own complete original proof of it (source0.tex lines 800--852, one \texttt{proof} environment of 53 lines). No mathematics is written, repaired or re-derived: statement and proof are the article's own text line for line. Required context retained uncounted: none. Every definition, display and label of those lines is retained unchanged. Omitted from this package and not redistributed: 1--782, 799--799, 853--1098. Nothing inside a retained excerpt is omitted or altered: every retained body is delivered exactly as the article's lines are written, so no omission receipt of any kind accompanies this package. Citations: the retained excerpts carry no citation command at all, so no bibliography entry of the article is transcribed and no reference is redistributed. Reference adaptation: none was needed and none was made; every reference inside the delivered unit resolves inside it, so no unresolved reference or citation is produced. Printed slips kept literal: no printed word, symbol, hypothesis or number of the counted statement or of its proof was altered. Layout and class: the article's own class is \texttt{amsart} with packages including hyperref, amsmath, verbatim, amsfonts, amssymb, blkarray, color, colortbl, xy, enumerate, geometry, mathrsfs, stmaryrd, bbold; the wrapper is the lane's pinned \texttt{amsart} editorial wrapper at 10pt on A4, which restates the theorem-like environments the retained text needs and adds the article's own macro definitions that the retained text needs (\texttt{\textbackslash geq}, \texttt{\textbackslash leq}), with extra wrapper packages (bm, bbm, dsfont, xspace, cancel, mathtools). The delivered numbering is the unit's own. Assets: no retained span contains a figure environment, a graphics-inclusion command, a PDF-page inclusion, a table or a TikZ picture; no image file is copied into this package, the package declares no asset dependency and no image file is redistributed with it. Licence: version arXiv:2509.17139v1 of this article is distributed under the Creative Commons Attribution 4.0 International licence, as recorded in the arXiv licence metadata for this exact version and shown on the abstract page https://arxiv.org/abs/2509.17139v1. A full rights notice is delivered at licenses/arxiv-2509.17139-CC-BY-4.0.txt. Characters: every retained excerpt is pure ASCII, so the delivered engine, XeLaTeX with its default Latin Modern text font, reports no missing character.
\begin{proposition}\label{prop:MM_4_4_1}
In the above situation the Herzog--Kunz sequence 
$\widetilde{a_1} < \cdots < \widetilde{a_N}$ of $S$ 
satisfies 
	$$ \{\widetilde{a_{i_0+1}},\ldots,\widetilde{a_{n+s}}\} 
    = \{a_{i_0+1}, \ldots, a_{n}, b_1,\ldots, b_s \} $$
and $S$ has a parametrization 
	$$ S = k[[\widetilde{x_1},\ldots,\widetilde{x_{n+s}}]] $$
with Herzog--Kunz generators $\widetilde{x_1},\ldots,
\widetilde{x_{n+s}}$  such that 
	$$ \begin{array} {l c l l}
	\widetilde{x_j} & = & x_j & \text{for }\, j \leq i_0 \\
	\widetilde{x_j} & = & t^{\widetilde{a_j}} & \text{for } 
    \, j > i_0 \\
	\end{array} $$  
\end{proposition}

\begin{proof}
The Herzog--Kunz sequence $\widetilde{a_1} < \cdots < 
\widetilde{a_N}$ of $S$ satifies $N = n+s$ as 
$\mathrm{edim}(S) = n+s$. Thus it suffices to show that 
	$$ \{ \widetilde{a_1} ,\ldots ,\widetilde{a_{n+s}}\} = 
	\{a_1, \ldots, a_n, b_1, \ldots b_s\} $$
for then for $j > i_0$, $\widetilde{a_j} \geq c_S$, and 
therefore $t^{\widetilde{a_j}} \in S$.

As $\mathfrak m_S = (x_1, \ldots, x_n, t^{b_1}, \ldots, 
t^{b_s})$, any $s \in \mathfrak m$ may be written as 
	$$ s = r + \sigma $$
with $r \in \mathfrak m$ and $v(\sigma) \geq b_1 $, and 
similarly any $s_2 \in \mathfrak m_S^2 = (x_i \cdot x_j, 
x_i \cdot t^{b_l}, t^{b_l} \cdot t^{b_m})$ 
may be written as 
	$$ s_2 = r_2 + \sigma_2 $$
with $r_2 \in \mathfrak m^2$ and $v(\sigma_2) \geq b_1 \,
(\geq c_R -a_1)$. Thus 
	$$ V_1 = \{w\in v(\mathfrak m)\,\vert\,w < c_R -a_1\} = 
	\{w\in v(\mathfrak m_S)\,\vert \, w < c_R -a_1\} = W_1 $$
and 
	$$ V_2=\{ w\in v(\mathfrak m^2)\,\vert \,w<c_R-a_1\} = 
	\{w \in v(\mathfrak m^2_S)\,\vert \,w < c_R -a_1\}=W_2 $$
implying that 
	$$ \{ \widetilde{a_i} \, \vert \, \widetilde{a_i} 
    \leq c_R-a_1 \} = W_1 \setminus W_2
	= V_1 \setminus V_2 =\{a_i\, \vert \,a_i\leq c_R-a_1 \} 
    = \{a_1, \ldots, a_{i_0} \} $$
It remains to show that 
	$$\{\widetilde{a_{i_0+1}},\ldots,\widetilde{a_{n+s}}\}= 
	\{ a_{i_0+1}, \ldots, a_n, b_1, \ldots, b_s \} $$
As both sets have the same cardinality, it suffices to show 
the inclusion $\subseteq$. 

Note that $\widetilde{a_{n+s}} < c_R$ as otherwise 
	$$ \widetilde{a_{n+s}} \in v(\mathfrak C_R) \subseteq 
    v(\mathfrak m^2) \subseteq v(\mathfrak m_S^2) $$
and similarly $a_n < c_R$. Thus both sets are contained in 
$\{c_R-a_1, \ldots, c_R-1\}$. 

If $u \in \{w \in V(S) \, \vert \, c_R-a_1 \leq w < c_R \} 
\setminus \{ a_{i_0+1}, \ldots, a_n, b_1, \ldots, b_s\}$, 
then $u \in V(R)$ (by the choice of $b_1, \ldots, b_n$) 
and thus 
$u \in v(\mathfrak m^2)$ (by the choice of $a_{i_0+1}, 
\ldots, a_n$), hence $u \in v(\mathfrak m_S^2)$, 
and therefore 
	$$ u \notin \{ \widetilde{a_{i_0+1}}, \ldots, 
    \widetilde{a_{n+s}} \} $$
completing the proof. 
 
\end{proof}

\end{document}
